\documentclass{article}\usepackage[T1]{fontenc}\begin{document}
For each fixed base object A, PayloadLaws.All.strong is a mathlib pseudonatural transformation from the actual weighted single-span payload action to the already-defined packed-iterate action. Its source index is the locally discrete bicategory of the ordinary category of Cat-valued presheaves and natural transformations. At Y its component is exactly the old evaluation equivalence's functor (strongComponent); its invertible constraint for T is precisely the symmetric orientation of the old evaluationPayload isomorphism (strongCell). The unit and composition equations hold at arbitrary whole-span objects, in particular length zero and every successor; the actions preserve identities and composition by equality of functors, while the evaluation comparison itself is only asserted invertible.

Read eq:telescope, lines 929--953 of the contextual-composition manuscript at 458ff36a4798e432dee5947afc402f9903a47f14. Distinguish three interfaces: the weighted single-span total category, the intermediate disjoint union of weighted length fibres, and the disjoint union of the iterated action. SumChange.evaluationPayload pastes the chosen inverse comparison with the tagged component comparison. Neither the base-variable unit/composition of MateLaws (\#201) nor the mixed hexagon of MixedMates (\#204) computes the identity or a pair of composable arrows in the payload variable.

PayloadLaws builds wholeAction, leafAction, and iterateAction as ordinary covariant functors at fixed A. fibreMap and allIterateMap supply the existing source/target maps and laws; sumId and sumComp account for transport through tags in the middle action. For the forward presentation from tagged weighted fibres to the whole weighted span, and for the tagged evaluation square, each comparison component is proved to arise from a genuine equality: its parallel transports agree by proof irrelevance. This argument does not apply to the inverse square of chosen equivalences. inverseUnit and inverseComp instead transfer the forward laws with MateLaws, retaining the units and counits of those equivalences. evaluationUnit and evaluationComp then paste across the same middle constraint.

The diagram versions state equality of natural transformations for the existing final comparison, not just a new witness with the right boundary. For T from Y to Z and U from Z to W the composite is T followed by U; the second isolated evaluation square is whiskered on the left by the whole action of T, the first on the right by the iterate action of U. ofDiagrams converts exactly those two laws to the unitors and associators of the strict bicategory Cat. Source 2-cell naturality reduces to equality in the locally discrete index.

evidence/payload-laws-author-fresh contains two complementary checks. In PayloadOrder, M adds a Boolean marker and T adds an option to an indexed family. On an empty input fibre the order M then T produces one element; T then M produces two. The file proves these output types are inequivalent and the two endotransformations unequal, computes an action on a nonidentity arrow, verifies that option misses a fibre arrow, and typechecks the final coherence at whole-span objects reached from tags zero and two. In PayloadMate, three chosen equivalences of a one-object integer-arrow category have the same underlying functors and different units. All forward faces are equality transports, but the two inverse faces and their composite have values -5, -10 and -15. Replacing the second face by an identity gives the wrong composite. This second check tests the generic mate argument; it does not purport to be a telescope instance.

Limits: split endpoints and the previous common-size base/payload regime are unchanged. At varying base, the base laws and mixed hexagon remain the results of \#201 and \#204. This tranche does not assemble them into a joint two-variable bicategorical equivalence, add arbitrary modifications between payload maps, or construct inverse transformation unit/counit modifications. It does not construct telescope flattening/multiplication, the monad comparison unit laws and associativity for flattening, or the algebra-recovery equivalence. \#105 remains open.
\end{document}
